
\section{Integrated Mission Simulation and DITL Supporting Analysis}\label{app:ditl-detail}
\subsection{Purpose and Source}\label{app:ditl-source}
The integrated model in \texttt{src/simulations/ditl/ANTARA\_Science\_Sim.py} combines orbit propagation, spacecraft attitude, electrical power, mode transitions, proton-environment assumptions and detector accumulation, and ground-station communications. Its current implementation is a simulation model, not a hardware-in-the-loop qualification system. The shared configuration and individual presets are documented in Appendix~\ref{app:sim-provenance}.
\subsection{Defined Mission Scenarios}\label{app:ditl-preset-details}
The source defines \texttt{leop\_24h}, \texttt{no\_leop\_24h} and \texttt{leop\_1yr}. The first includes Cold and Dark and Detumble; the second begins in Sun Pointing without LEOP startup; the annual scenario includes startup over 8,760 hours. Command-line values may override preset defaults. A comparison of the three scenarios is meaningful only when each output is accompanied by its exact version and configuration.
\subsection{Mode Execution, Time and Energy}\label{app:ditl-accounting}
During startup, detumble and slewing, the simulator uses its fine time cadence. Other operating modes may use a coarser step. Mode occupancy, energy use, detector accumulation and mission elapsed time advance with the same actual integration interval. The interval length is recorded in \texttt{interval\_duration\_sec}; the last zero-duration terminal telemetry row should not be counted as additional mission time. Reported fractions should therefore use total elapsed time, not raw row fractions.
The ADCS model retains finite slew and settling requirements, momentum saturation handling and recovery logic. During saturation hold, scientific collection and downlink transmission are inhibited in the documented model. The end of a planned slew duration is not by itself proof of achieved target alignment.
\subsection{Science Data, Detector Statistics and Model Limits}\label{app:ditl-science}
The detector model accumulates expected bin counts and serializes histograms into science data separately. An interval with accepted events need not produce a positive new packet volume until a flush. SAA-exit and simulation-horizon flushes belong to the corresponding pass totals. Bin precision estimates use $1/\sqrt{N}$ for accumulated expected counts, but this does not represent calibrated detector response, stochastic-count realization or a fully validated low-count confidence interval. The environment model's legacy AP9-named variable is a configured baseline and not an independent AP9 computation.
\subsection{Results and Acceptance Evidence}\label{app:ditl-evidence}
The source archive contains figures with scenario names but the final report must associate each approved graphic with its source revision, scenario configuration, complete CSV and summary. Earlier figures produced before corrected step accounting, attitude or wheel-momentum handling should not be mixed with current simulations without review. The code and documentation alone establish model capabilities and algorithms, not full-year completion, numerical accuracy or mission success.
\begin{table}[htbp]\centering\caption{Integrated DITL results required for the final evidence package.}\label{tab:ditl-closure}\begin{tabularx}{\linewidth}{p{0.25\linewidth}XX}\toprule Scenario & Main quantities to extract & Evidence status\\\midrule
LEOP 24 hours & Start-up and detumble times, wheel behavior, minimum battery, first contact, SAFE events & Run-specific review required\\
Nominal 24 hours & Mode occupancy, SAA active time, packet accumulation, downlink, pointing error, power and battery extrema & Run-specific review required\\
Integrated one year & Year-end battery and orbit, total science and contact opportunities, cycle and anomaly statistics, run completion & Full-run verification required\\\bottomrule\end{tabularx}\end{table}
